\section{Method}
\label{sec:method}
In this section, we first introduce the framework of our proposed \textbf{\textit{RoboOS}} and explain the functionalities of its three key components. Next, we outline the primary workflow pipeline for multi-robot collaboration within \textbf{\textit{RoboOS}}, detailing the implementation of hierarchical information interaction. 
Finally, we present optimizations for edge-cloud communication and cloud-based distributed inference, ensuring high-frequency interactions and enabling large-scale deployment.

\begin{figure*}[t]
    \centering
   \vspace{-0.2cm}
    \includegraphics[width=0.97\linewidth]{images/overview.pdf}
    \caption{
    \textbf{Framework of RoboOS.} RoboOS is a Brain-Cerebellum hierarchical architecture for multi-robot collaboration that comprises three core components: a) Cloud-based Embodied Brain model for high-level cognition and multi-agent coordination; b) Distributed Cerebellum Modules for executing robot-specific skills; and c) Real-Time Shared Memory for enhanced environmental awareness.
        }
       \vspace{-0.5cm}
    \label{fig:method}
\end{figure*}

\subsection{Framework of RoboOS}
\label{subsec:framework}
As shown in Fig.~\ref{fig:method}, \textbf{\textit{RoboOS}} is a unified embodied system built on a biologically inspired Brain-Cerebellum hierarchical architecture, comprising three core components: the Embodied Brain Model (\texttt{RoboBrain~\cite{ji2025robobrain}}), the Cerebellum Skill Library, and Real-Time Shared Memory. Deployed via the FlagScale MLLMs toolkit~\cite{FlagScale2024}, the edge-cloud RoboOS framework facilitates seamless multi-robot coordination by synchronizing cognition across the multiple agents.
The system operates as follows: First, the Embodied Brain Model manages system-wide tasks, including multi-robot task planning, tool invocation, spatiotemporal memory updates, and adaptive error correction through continuous three-level feedback loops. Second, the Cerebellum Skill Library, deployed on individual robot terminals, offers modular, plug-and-play functionality via standardized Robot Profiles. Finally, the Redis-optimized Shared Memory maintains a dynamic knowledge base of spatial relationships, operational states, and historical data to support real-time decision making. This architecture ensures robust large-scale deployment while maintaining the low-latency interactions essential for embodied AI systems.

\textbf{Embodied Brain Model (RoboBrain)} 
The cloud-deployed MLLM can be any existing model, including proprietary models~\cite{hurst2024gpt4o, anthropic2024claude3.5, Google2023gemini} or open-source options~\cite{liu2024deepseek,Qwen2.5-VL,chen2024internvl,li2024llavaov,abdin2024phi,beyer2024paligemma}. To better suit embodied scenarios, we adopt RoboBrain~\cite{ji2025robobrain} as our Embodied Brain Model, enhancing its capabilities for the RoboOS framework. Building on RoboBrain's functionalities—single-robot planning, affordance prediction, and trajectory prediction—we implement multi-stage training to improve multi-robot task planning, agent-based tool invocation, and spatiotemporal updates using the pretrained Qwen2.5VL-7B model~\cite{Qwen2.5-VL}. Key enhancements include \textit{Multi-Robot Task Planning}, utilizing real-time shared spatiotemporal memory to predict workflow topologies for collaborative tasks; \textit{Agent-Based Tool Invocation}, managing agents and invoking tools as needed with self-corrective planning based on feedback; \textit{Spatiotemporal Memory Update}, dynamically updating shared memory in real time according to subtask execution and tool feedback; and \textit{Low-Level Guidance}, predicting manipulable regions and trajectories during tool execution to assist in manipulation~\cite{ji2025robobrain}.

\textbf{Cerebellum Skill Library} 
This modular, plug-and-play embodied toolkit supports various robotic embodiments (\textit{e.g.}, single-arm, dual-arm, wheeled, humanoid) with low-latency execution for manipulation and navigation throughout the robotic task cycle. 
The Cerebellum Skill Library encompasses three key aspects: manipulation types, integrating both expert-based tools (\textit{e.g.}, affordance-aware grasping~\cite{tang2025affordgrasp}, generalized grasping~\cite{fang2023anygrasp}) and VLA-based tools (\textit{e.g.}, OpenVLA~\cite{kim2024openvla}, RDT-1B~\cite{liu2024rdt}, Octo~\cite{team2024octo}, $\pi_0$~\cite{black2024pi_0}); 
navigation types, supporting traditional mapping-localization-navigation pipelines (\textit{e.g.}, SLAM~\cite{xu2025fastliolocalization, grisetti2010tutorial}) and vision-language-navigation (VLN) tools (\textit{e.g.}, MapNav~\cite{zhang2025mapnav}, NavID~\cite{zhang2024navid}); 
and specialized skills for contact-rich interactions, deformable object handling, and dexterous hand control~\cite{fu2025cordvip, zhong2025dexgrasp, aksoy2025planning}. 
Standardized tool and robot profiles ensure seamless integration and interoperability across diverse robotic platforms.

\textbf{Real-Time Shared Memory} 
%This component maintains spatial, temporal, and robotic memory to enable fault prediction and load balancing across robots. 
%\textit{Spatial memory} consists of a scene graph~\cite{chang2021comprehensive} that tracks real-time spatial relationships, object locations, and robot positions. 
%\textit{Temporal memory} records task execution history, feedback, tool-calling logs, and other temporal data to support adaptive decision-making. \textit{Robotic memory} stores real-time system attributes such as motion domain constraints, joint states, and battery levels, optimizing task allocation based on each robot's capabilities, power status, and connectivity conditions.
This component maintains spatial, temporal, and robotic memory to enable robust coordination and adaptive decision-making.  
\textit{Spatial memory} is represented as a dynamic scene graph~\cite{chang2021comprehensive}, structured hierarchically into floors, rooms, and object-level nodes~\cite{werby2024hovsg,gu2024conceptgraphs}. Each node encodes attributes such as position, affordances, and semantic labels, while edges represent spatial or functional relationships. Scene graphs are constructed from multi-view RGB-D inputs using SAM segmentation and CLIP features, merged via geometric clustering and open-vocabulary grounding.  
\textit{Temporal memory} records task execution history, feedback, tool-calling logs, and other temporal data to support adaptive decision-making.
\textit{Robotic memory} stores real-time system attributes such as motion domain constraints, joint states, and battery levels, optimizing task allocation based on each robot's capabilities, power status, and connectivity conditions.



\begin{figure*}[t]
    \centering
    \vspace{-0.2cm}
    \includegraphics[width=0.97\linewidth]{images/pipeline.pdf}
    \caption{
    \textbf{Pipeline of RoboOS.} The RoboOS framework implements a workflow pipeline for multi-robot collaboration, consisting of four key phases: (1) hierarchical task decomposition, (2) topology-aware subtask allocation, (3) distributed agent-based execution, and (4) dynamic memory updating. This integrated workflow enables coordinated task completion while maintaining adaptability to environmental and operational constraints.
    }
       \vspace{-0.2cm}
    \label{fig:pipeline}
\end{figure*}

\subsection{Workflow Pipeline of RoboOS}
\label{subsec:pipeline}
The proposed RoboOS demonstrates high task concurrency and flexibility in multi-robot task allocation. To clarify the overall workflow pipeline of RoboOS, we use a single global task for detailed elaboration, as shown in Fig. \ref{fig:pipeline}.

\textbf{Step 1: Global Task Decomposition} Upon receiving the global task instruction \( T_{\text{global}} \), RoboOS initiates a Retrieval-Augmented Generation (RAG) process via RoboBrain to query the shared spatial memory, extracting environment-relevant information \( M_s \). This is integrated with (i) state feedback \( M_t \) from prior task executions (stored in shared temporal memory), (ii) the robot’s operational status \( S_r \) (\texttt{idle}, \texttt{busy}, or \texttt{offline}), (iii) the robot skill library \( M_r \), and (iv) \( T_{\text{global}} \). RoboBrain processes these inputs to generate a structured reasoning trace \( \mathcal{R} \) and a subtask graph \( \mathcal{G} \), formalized as: 
\begin{equation}
    (\mathcal{R}, \mathcal{G}) = \text{RoboBrain}\big(M_s \oplus M_t \oplus S_r \oplus M_r \oplus T_{\text{global}}\big),
\end{equation}
where \( \oplus \) denotes the concatenation or fusion of multimodal inputs.

\textbf{Step 2: Topological Subtask Allocation} 
The Monitor dynamically schedules and allocates subtasks in parallel based on the topological dependencies encoded in the directed acyclic graph \( \mathcal{G} \). 
Each subtask in \( \mathcal{G} \) is classified into two types: \textit{(1) Single-Robot Subtask \((d, r_i)\)}, executed autonomously by robot \( r_i \) at topological depth \( d \); and \textit{(2) Collaboration Subtask \((d, r_{i:j})\)}, requiring coordinated execution among multiple robots \(\{r_i, \dots, r_j\}\) at depth \( d \). Here, \( d \) represents the execution priority, while \( r_i \) (or \( r_{i:j} \)) denotes the assigned robot(s). 
To enforce dependency constraints, the Monitor employs \textit{\textbf{Parallel Allocation}}—executing independent subtasks concurrently at the same depth (\textit{e.g.}, \((1, R_1)\) and \((1, R_2)\) in Fig. \ref{fig:pipeline})—and \textit{\textbf{Sequential Allocation}}, where subtask \((d+1, r_k)\) is blocked until all prerequisites at depth \( d \) are fulfilled (\textit{e.g.}, \((2, R_{1+2})\)). 
In practice, the system supports concurrent management of multiple subtask graphs \(\{ \mathcal{G}_1, \mathcal{G}_2, \dots, \mathcal{G}_n \}\), ensuring real-time adaptability to dynamic robot states and evolving task dependencies.

\textbf{Step 3: Distributed Subtask Agent}  
For each subtask, RoboOS deploys a dedicated \textit{\textbf{Robotic Agent}} to manage execution. The Agent autonomously orchestrates tool selection from the Skill Library based on: (1) feedback from prior executions, (2) tool-calling history, and (3) partial spatial memory of the environment. This closed-loop reasoning facilitates dynamic error recovery.  
For example (Fig. \ref{fig:pipeline}), when tasked with ``Search for an egg and place it on the table'', the Agent sequentially invokes tools (\textit{e.g., ``detect an egg''}). If the search fails (\textit{e.g., no egg detected in the kitchen}), the Agent uses spatial memory to infer potential locations (\textit{e.g., the fridge}) and selects the navigation tool to ``move to fridge'', showcasing adaptive recovery through iterative tool refinement.

\textbf{Step 4: Dynamic Memory Updating}  
%Upon completing a subtask (whether successful or not), the shared memory is updated. For example, if the subtask ``Search for an egg and place it on the table'' succeeds, the egg's location is updated from ``kitchen'' to ``table'' via RoboBrain, which generates instructions to modify the Spatial Memory. Additionally, feedback, tool-calling history, and robot states are logged in the Temporal Memory and Robotic Memory.
Scene graphs are updated incrementally as robots perceive and act. New observations or interaction effects (e.g., filling, moving) modify node attributes and spatial edges in real time. For instance, after placing a cup on a shelf, its position and supporting surface edge are updated. Occlusion and history-aware correction ensure continuity, even when objects are temporarily hidden.  
Additionally, feedback, tool-calling history, and robot states are logged in the Temporal Memory and Robotic Memory.
% Temporal and robotic memories are likewise refreshed: each subtask updates execution logs, robot states, and tool calls, enabling persistent, cross-task learning in long-horizon scenarios.

\subsection{Edge-Cloud Deployment}
\label{subsec:edge-cloud}
Built upon our parallel training and inference framework, FlagScale \cite{FlagScale2024}, RoboOS supports end-cloud collaboration for multi-robot systems, creating a unified foundation for embodied intelligence. Designed for ``multi-robot, multi-modal, multi-task'' scenarios, it offers exceptional scalability and ultra-low-latency responsiveness. In edge deployments, robots automatically establish bidirectional communication with the cloud-based RoboBrain upon registration, enabling real-time task scheduling and status feedback via an efficient publish-subscribe mechanism (average command response latency $< 0.001 \, \text{s}$).
To manage the vast perception and behavioral data generated during long-term operations, FlagScale includes a memory-optimized data access engine that supports TB-level historical data with in-memory random access. This facilitates task replay, anomaly backtracking, cross-task knowledge transfer, and other critical scenarios. Additionally, FlagScale's framework supports parallel inference and multi-task cooperative scheduling of large models across distributed devices, unlocking the systemic potential of RoboBrain.